\def\stat{somin}



\def\tit{ИСПОЛЬЗОВАНИЕ ХЕШ-ФУНКЦИЙ ДЛЯ~ПОВЫШЕНИЯ

СКОРОСТИ МОРФОЛОГИЧЕСКОГО АНАЛИЗА

РУССКИХ ТЕКСТОВ}



\def\titkol{Использование хеш-функций для~повышения

скорости  морфологического анализа}

%русских текстов}



\def\autkol{Н.\,В.~Сомин, М.\,М.~Шарнин}



\def\aut{Н.\,В.~Сомин$^1$, М.\,М.~Шарнин$^2$}



\titel{\tit}{\aut}{\autkol}{\titkol}



%{\renewcommand{\thefootnote}{\fnsymbol{footnote}}

%\footnotetext[1]

%{Работа выполнена при частичной финансовой поддержке  программы

%<<Интеллектуальные информационные технологии, системный анализ и

%автоматизация>> (проект 1.7).}}







\renewcommand{\thefootnote}{\arabic{footnote}}

\footnotetext[1]{Институт проблем информатики Российской академии наук, somin@post.ru}

\footnotetext[2]{Институт проблем

информатики Российской академии наук, mc@keywen.com}





     \label{st\stat}



                  \thispagestyle{headings}



    \Abst{Рассматривается проблема повышения эффектив\-ности

морфологического анализа русского языка. Описывается разработанная в

ИПИ РАН система морфологического анализа: набор морфологических

характеристик и алгоритмы работы. Указываются задачи и сис\-те\-мы,

связанные с проб\-ле\-мой ло\-ги\-ко-се\-ман\-ти\-че\-ско\-го анализа текстов, в

которых данная морфологическая сис\-те\-ма нашла применение. Обсуждаются

особенности сис\-те\-мы с точки зрения занимаемой памяти и ско\-рости работы.

    Предлагается способ хранения мор\-фо-лек\-си\-че\-ской информации с

помощью хеш-функ\-ций, обеспечивающих высокую ско\-рость доступа.

Обсуждаются трудности, возникающие при реализации такого подхода, и

рассматриваются пути их преодоления. Приводится структура

информационных массивов новой версии и реализованные в ней поисковые

алгоритмы, а также даются сведения по подсистеме ввода и корректировки

морфологической информации. Приводятся конкретные параметры новой

реализации морфологической сис\-те\-мы и данные по ускорению работы по

сравнению с предыдущей версией.

    В заключение обсуждаются возможности по развитию новой версии

морфологии и перенесению предложенного подхода к реализации на другие

компоненты лингвистического процессора.}



    \KW{морфологический анализ; лингвистический процессор;

морфологическая омонимия; хеш-функция}



\DOI{10.14357\!/\!08696527140315}







\vskip 14pt plus 9pt minus 6pt





      \thispagestyle{headings}





    \section{Введение}



    Задача морфологического анализатора~--- нормализация слов, наиболее

полное и точное определение морфологических признаков лексем, а также (в

ряде случаев) нахождение их семантического класса. Отметим, что к

настоящему времени разработан целый ряд морфологических анализаторов

русского языка, среди которых можно отметить~[1--3]. Авторы принимали

участие в создании разработанной в ИПИ РАН системы морфологического

анализа~[4], а~также ее интеграции в состав целого ряда лингвистических

процессоров: <<Криминал>>, <<Аналитик>>, <<Поток>>,

<<Лингвопроцессор>>, <<Лингво-ИИ>>~[5--11] и~ряда других. Помимо

этого система морфологического анализа применялась при реализации

экспериментальных лингвистических процессоров ЭССЕИСТ и~\mbox{ТЕРМИН-3}~[12], а также в системах рубрицирования текстов~[13, 14]. Опыт

эксплуатации морфологической системы отражен в~работе~[4]. Везде

входной поток текстов, проходя через морфологический блок

(интегрированный с блоком лексического анализа), дополнялся подробной

морфо-лексической информацией, на основе которой уже производился

дальнейший анализ.



    Однако в связи с доступностью на современном этапе огромного

количества текстов из Интернета и, соответственно, необходимостью

обработки больших объемов данных (Big Data) к системе морфологического

анализа предъявляются новые требования. В~частности~--- требование

высокой производительности. Причем, как представляется, уповать на

простое увеличение скорости работы современных процессоров нет

оснований, поскольку темпы роста их вычислительной мощности явно

отстают от темпов увеличения объемов текстов, требующих обработки.

Поэтому необходимо искать другие пути повышения эффективности

морфологического анализа. Они были найдены в использовании ресурса

оперативной памяти современных компьютеров, который за последнее

десятилетие вырос более чем на порядок. Именно за счет использования

большего объема памяти в новой версии морфологической системы удалось

существенно повысить скорость анализа, не потеряв (и даже увеличив)

функциональные возможности старой версии системы.



\section{Возможности системы морфологического анализа}





    Разработанная система морфологического анализа позволяет

обрабатывать текст на русском и английском языках, выполняя для каждого

слова определение его канонической формы и выдавая его полный

морфологический анализ. Кроме того, система позволяет:

    \begin{itemize}

\item выполнять генерирование всей парадигмы словоформ по заданной

слово\-форме;

\item формировать предметные терминологические словари и

идентифицировать термины с учетом парадигматических изменений~[15].

    \end{itemize}



    Ниже представлен набор морфологических признаков с указанием

соответствующих символов (на первом месте) и формируемых на их основе

фраг\-мен\-тов-при\-зна\-ков (в скобках~--- они используются программами

семантического анализа на языке ДЕКЛ~\cite{8-som}):



\begin{tabular}{rcl}

    f &---& фамилия (FAM);\\

    i&---& имя (NAME\_);\\

    h&---& отчество (NAME\_2);\\

    N&---& существительное (NOUN);\\

    A&---& прилагательное (ADJ);

    \end{tabular}



\noindent

    \begin{tabular}{rcl}

    . &---& единственное число (по умолчанию);\\

    :&---& множественное число (MANY);\\

    м&---& мужской род;\\

    ж&---& женский род;\\

    с&---&средний род;\\

    и&---& именительный;\\

    р&---& родительный (КОГО) ;\\

    д&---& дательный (КОМУ);\\

    в&---& винительный;\\

    т&---& творительный (КЕМ);\\

    п&---& предложный;\\

    D&---& деепричастие (DPRICH);\\

    R &---& наречие (ADJ);\\

    G&---& междометие (MZ);\\

    U&---& числительное (NUM\_W);\\

    O&---& отглагольное прилагательное типа <<\textit{бывший}>>,

<<\textit{битый}>> (ADJ\_O);\\

    H&---& частица (PART);\\

    P&---& предлог (PREP);\\

    M&---& местоимение (PRON);\\

    T&---& причастие (PRICH);\\

    V&---& глагол (VERB);\\

    S&---& союз (CONJ);\\

    E&---& модальное слово (VERB\_2);\\

    Z&---& указательное местоимение <<\textit{тот}>>, <<\textit{этот}>>

(PRON);\\

    W&---& вводное слово типа <<\textit{надо}>>, \ldots  (AUX);\\

    =&---& настоящее;\\

    $>$&---& будущее;\\

    $<$&---& прошедшее;\\

    +&---& положительная степень (по умолчанию);\\

    $^*$&---& сравнительная степень;\\

    \$&---& превосходная степень;\\

    л&---& 1-е лицо;\\

    Л&---& 2-е лицо;\\

    L&---& 3-е лицо;\\

    П&---& полная форма (по умолчанию~--- FULL);\\

    к&---& краткая форма (SHORT);\\

    Г&---& одновременно существительное\;+\;при\-ла\-га\-тель\-ное,

числительное,\\

&& причастие (<<\textit{малый}>>, <<\textit{нищий}>>);\\

    ф&---& инфинитив (INFIN);\\

    ъ&---& изъявительное наклонение;\\

    !&---& повелительное наклонение;\\

    я&---& форма без -\textit{ся}/-\textit{сь} (по умолчанию);\\

    I&---& одушевленное;

    \end{tabular}



\noindent

    \begin{tabular}{rcl}

        X&---& неодушевленное;\\

    )&---& безличный глагол;\\

    \{ &---& личный/безличный глагол;\\

    $[$&---& личное местоимение;\\

    v&---& личный глагол;\\

    \}&---& личный/безличный глагол;\\

    /&---& совершенный/несовершенный вид;\\

    ш&---& совершенный вид;\\

    н&---& несовершенный вид;\\

    -&---& отрицательное местоимение;\\

    F&---& определительное местоимение;\\

    Ж&---& притяжательное местоимение;\\

    q&---& порядковое числительное;\\

    Q&---& количественное числительное;\\

    ?&---& относительное (вопросительное) местоимение;\\

    s&---& собирательное местоимение;\\

    ч&---& непереходный глагол;\\

    $\sim$&---& переходный глагол;\\

    @&---& переходный/непереходный глагол;\\

    w&---& слово в значении сказуемого (например, <<\textit{цап}>>);\\

    u&---& неопределенное местоимение (<<\textit{некоторый}>>);\\

    v&---& возвратный глагол (с \textit{ся}/\textit{сь});\\

    d&---& действительный залог;\\

    b&---& страдательный залог;\\

    ю&---& разносклоняемое;\\

    Б&---& архаичное;\\

    Д&---& числительное, характеризуется числом;\\

    И&---& частица без ударения;\\

    Й&---& ударный предлог;\\

    Н&---& 1-е склонение;\\

    Ф&---& 2-е склонение;\\

    Ц&---& 3-е склонение;\\

    3&---& группа местоимения;\\

    Ч&---& 1-е спряжение;\\

    Ш&---& 2-е спряжение;\\

    ы&---& инициалы (NAME!);\\

    6&---& в слове более одной прописной (HEAD\_1);\\

    7&---& слово начинается с прописной (NAME0);\\

    0&---& в слове все прописные (HEAD);\\

    э&---& английское слово (ENGL);\\

    8&---& целое число (NUM);\\

    9&---& число с точкой (NUM\_F);\\

    щ&---& сложный предлог;

        \end{tabular}



\noindent

    \begin{tabular}{rcl}

        t&---& форма английского глагола \textit{be};\\

    R&---& артикль;\\

    Я&---& личное местоимение (PRON\_1);\\

    \#&---& морфология определена по аналогии;\\

    а&---& англ.\ \textit{ing}-форма (WEB\_C);\\

    б&---& англ.\ \textit{ed}-форма past-1 (MAIL\_E);\\

    Е&---& англ.\ \textit{ed}-форма past-2.

\end{tabular}



\smallskip





    Для одной лексемы блок морфологического анализа выдает множество

признаков, набор которых и характеризует морфологический тип. Кроме

морфологических блок выдает еще множество лексических признаков, о

которых речь шла выше, а также ряд фонетических, которые могут быть

использованы для синтеза речи.



    Особое место занимает признак <<\#>>. Он означает, что данный набор

признаков сформирован <<по аналогии>>, т.\,е.\ была найдена словоформа с

окончанием (т.\,е.\ несколькими конечными буквами), таким же как у данной

лексемы, и~морфологические признаки этой словоформы были приписаны

данной лексеме. Варианты разбора <<по аналогии>> применяются для

лексем, которых нет в~морфологическом словаре.



\section{Борьба с морфологической омонимией}



    Использование обобщенного словаря основ (ОСО) может приводить к

лишним вариантам морфологического анализа. Было предложено два

эффективных способа борьбы с морфологической омонимией.



    Первый способ~--- эмпирический алгоритм, отбрасывающий наименее

вероятные варианты морфологического анализа. Такая <<зачистка>>

вариантов выполняется по многим критериям, учитывающим наличие слова

в словаре основ (СО), длину основы со словарем хвос\-тов основ (СХО), часть речи. Эмпирический алгоритм фиксирует

все варианты разбора в порядке их вероятности, что необходимо для ряда

приложений, когда используется только один вариант морфологического

анализа.



    Второй способ~--- частичный синтаксический анализ. В~предложении

слово вступает в синтаксические связи с другими словами, и выявление этих

связей позволяет отбросить варианты морфологического анализа, этим

связям не удовлетворяющие. Прежде всего было реализовано распознавание

двух конструкций: полного согласования и генетической цепочки.



    Алгоритмы борьбы с морфологической омонимией являются основным

содержанием так называемого процесса <<постморфологии>>, т.\,е.\ обработки

результатов морфологического анализа на завершающем этапе. Как показала

практика, эти процессы занимают в морфологической системе большое

место~--- как по алгоритмической сложности, так и по времени их

выполнения.



\section{Особенности блока английской морфологии}



    Английский морфологический словарь использует разработанное для

русской морфологии программное обеспечение, которое оказалось

возможным адаптировать к специфике английского языка. Блоки английской

и русской морфологии выдают практически одни и те же морфологические

характеристики, что позволяет использовать для синтаксико-семантического

анализа англоязычных текстов те же средства, что и для русского языка.

В~результате появилась возможность записывать лингвистические знания

для этих языков в одном и том же формализме.



    Общий объем словаря основ блока английской морфологии~--- около

85~тыс.\ слов. Для повышения качества работы этого блока в него был

введен ряд специфических для английского языка алгоритмов, которые в

основном касаются отсева лишних вариантов морфологического анализа,

поскольку слова английского языка чрезвычайно омонимичны~--- очень

часто одно и то же слово может быть и существительным, и глаголом, и

прилагательным. В~блоке английской морфологии реализованы алгоритмы,

позволяющие корректно отбрасывать лишние варианты (другие варианты

отсеиваются в процессе син\-так\-ти\-ко-се\-ман\-ти\-че\-ско\-го анализа).

Блок лексико-морфологического анализа был модифицирован для работы с предметными словарями

английского языка, которые удалось совместить со словарями русского

языка.



\section{Анализ существующей схемы разбора}



    Реализованная ранее схема морфологического анализа выглядела

следующим образом.



    Сформированы словарь классов окончаний (СКО)~--- все возможные

парадигмы русского языка и СО~--- основы слов со

ссылками на соответствующий класс окончаний. В~основе хранится

неизменяемая часть слова с постоянными для данного слова

морфологическими признаками. Класс окончаний содержит все окончания

данной парадигмы плюс изменяемые (т.\,е.\ присущие каждому окончанию)

морфологические параметры.



    Например, слово <<\textit{бытие}>> имеет основу <<\textit{быти}>> и

класс окончаний за номером~1759, содержащий окончания в именительном,

родительном, дательном, винительном, творительном и предложном

падежах, а именно: <<\textit{е}>>, <<\textit{я}>>, <<\textit{ю}>>,

<<\textit{е}>>, <<\textit{ем}>>, <<\textit{и}>> (множественного числа это

слово не имеет). Соответственно в~СО имеется запись <<\textit{быти

1759}>>, а в СКО под номером~1759 закодирована парадигма с указанными

окончаниями.



    Легко видеть, что в общем случае может иметь место омонимия основ,

т.\,е.\ несколько одинаковых основ, но с разными парадигмами окончаний.

С~другой стороны, одна и та же парадигма может <<обслуживать>> много

разных основ. Именно поэтому разделение на СО и СКО

оказывается с точки зрения экономии пространства столь

эффективным приемом. Однако поиск всех корректных вариантов разбора

требует значительного времени. Действительно, необходимо проверить все

варианты разбиения входной лексемы на основу и окончание и~для каждого

варианта найти в~словаре основу и~проверить, имеется ли во всех классах

окончаний, связанных с~этой основой, предполагаемое окончание.



    В предыдущей версии был также реализован алгоритм обработки

незнакомых системе слов <<по аналогии>>~[4]. В~реализации этого метода

использовался третий словарь~--- СХО,

в который записываются все 1-, 2-, 3-бук\-вен\-ные

и~т.\,д.\ <<хвос\-ты>> основ (первые буквы основ отбрасываются) с указанием

соответствующего класса окончаний. Было решено, что в СХО не будет

одинаковых <<хвос\-тов>>, а его класс окончаний вычисляется из

статистических соображений~--- по максимуму основ в СО, имеющих

данный <<хвост>> и данный класс окончаний. Если слово не находится в

 СО, то та же схема анализа повторяется, но уже с~по\-мощью пары

словарей СХО--СКО.



    Отметим, что в предыдущей версии все словари были реализованы в

виде деревьев, каждый уровень которых соответствовал одной букве (причем

порядок в основах и окончаниях был инвертированным). Это давало

возможность резко уменьшить рабочий объем словарей. Однако поиск

осуществлялся далеко не лучшим в смысле времени работы образом: путем

поиска нужной ветки и спуска на следующий уровень по веткам дерева.



    Описанную реализацию можно было улучшать, добиваясь

определенного прироста быстродействия. Но достичь качественного

улучшения представлялось проблематичным.



\section{Принципы реализации в новой версии}



    Прежде всего ставилась задача кардинально повысить скорость анализа

русских текстов. Оказалось, что это можно осуществить только путем отказа

от принципа разделения на СО и СКО и

объединения всей морфологической информации в единый словарь.

Разумеется, это резко повышает требования к оперативной памяти. Но в

современных компьютерах она уже достигает нескольких гигабайт, что

позволяет реализовать этот принцип даже для словаря порядка 100~тыс.\

слов.



    Таким образом, схема информации представляется нарочито

упрощенной: единый словарь всех словоформ с полным морфологическим

описанием каждой словоформы. Каждая словоформа рассматривается как

ключ доступа, по которому производится прямое чтение информации и

извлечение морфологических знаний.



    Однако, чтобы такая <<примитивная>> схема реализации была

успешной, требуется решить несколько принципиальных задач. Прежде

всего, это задача прямого доступа по ключу для больших массивов.

В~данной реализации предлагается использование принципа хеш-функ\-ций,

позволяющего реализовать очень быстрый доступ. Принцип работы

    хеш-функ\-ций хорошо известен. Вкратце он сводится к преобразованию

ключа в уникальный адрес доступа с~по\-мощью ряда математических и

логических операций. Хотя такие вычисления занимают мало времени,

однако при реализации этого подхода возникает ряд нетривиальных проблем.

Во-пер\-вых, это достижение равномерности такого преобразования

(в~смысле равновероятного покрытия всего адресуемого пространства).

Указанная равномерность так или иначе решается разработчиками алгоритма

хеширования. Во-вто\-рых, это проблема <<коллизий>>, т.\,е.\ разрешения

конфликтных ситуаций, когда несколько разных ключей преобразуются в

один и тот же адрес. И,~наконец, в-третьих, это проблема рационального

заполнения <<пус\-тот>>, т.\,е.\ элементов массива, остающихся пустыми в

силу того, что в них не отображается ни один ключ. Вторые две проблемы

необходимо уже решать разработчику системы, использующей

    хеш-функ\-ции.



    Помимо корректного использования хеш-функ\-ций имеются еще

обычные для такого рода систем проблемы ввода и корректировки

морфологической информации. Ниже будут описаны методы, позволяющие

это проблемы удовлетворительно решить.



\section{Структура информационных массивов.}



    Ради рационального использования памяти вводятся два массива

информации: массив адресации (МА) и массив информации (МИ).



    Массив адресации представляет собой хеш-массив, в котором располагаются

ссылки (смещения) на МА. Этот массив организован в зоны по несколько

записей в зоне, причем в 1-м байте зоны хранится число записей.

    Хеш-функ\-ция преобразует ключ (словоформу) в два хеш-кода: Hash1,

который указывает на номер зоны, и Hash2~--- 3-байт\-ный код, которым

идентифицируется запись внутри зоны. Каждая запись включает в себя поле

Hash2 и адрес в МИ. Все записи, попадающие в одну и ту же зону (т.\,е.\

имеющие один и тот же Hash1), последовательно заполняют эту зону и

ищутся в ней по Hash2. При переполнении зоны они пишутся в следующую

зону на свободное место.



    Двукратное хеширование позволяет достаточно экономно использовать

адресуемое пространство, практически не теряя в скорости доступа. Массив

адресации управляется несколькими параметрами: величиной массива (в байтах),

длиной записи и длиной зоны. Это дает возможность после нескольких

попыток подобрать параметры, оптимально подходящие под величину

массива морфологической информации.



    Массив информации является текстовым массивом, в котором через

    вы\-чис\-ли\-тельный комплекс хранятся

(в~определенном формате) словоформы и морфологические признаки.

Информация в нем структурируется по меточному принципу: информация

делится на части, каждая из которых предваряется метками, такими как

<<m=>>~--- меткой морфологических параметров, <<k=>>~--- меткой

семантической информации и~т.\,д. Например:



    высадки =$>$ $\vert$ ка m=52NжXНр $\vert$ка m=52NжXНви: $\vert$ок

m=52NмXФви:



    Здесь описан морфологический разбор словоформы <<высадки>>,

причем дается три варианта разбора: родительный падеж единственного числа слова

<<высадка>>; множественное число того же слова и множественное число

существительного <<высадок>>. После символа <<$\vert$>> дается

окончание канонической формы, а две цифры после <<m=>> означают длину

основы и длину окончания для разных вариантов разбора слова

<<высадки>>.



\section{Анализ <<по аналогии>> в новой версии}



    В принципе, <<прямой>> подход к построению морфологического

словаря может быть применен и для анализа <<по аналогии>>. Для этого

необходимо сгенерировать все усеченные варианты каждой словоформы (с

указанием предположительных морфологических признаков), а в процессе

обработки входной словоформы, при отсутствии ее в словаре, начать процесс

усечения с начала слова с поиском каждого усеченного варианта.



    Однако эксперименты показали, что при этом словарь (точнее оба

массива: МА и МИ) увеличиваются более чем на порядок, что уже вызывает

затруднение при реализации даже на современных персональных

компьютерах (тем более, что морфологический анализ~--- лишь одна из

многочисленных задач, решаемых в рамках лингвистического процессора).

Поэтому было решено реализовать гибридный подход: слова из словаря

распознавать с помощью новой версии, а~для анализа <<по аналогии>>

оставить алгоритмы предыдущей версии.



\section{Корректировка и ввод информации}



    Ввод в МА и МИ осуществляется с помощью специального

блока загрузки с морфологического массива (ММ). Этот массив имеет вид

текстового файла и может корректироваться с помощью любого текстового

редактора. Порядок записей в этом файле не имеет значения. Однако

требуется, чтобы все варианты морфологического разбора одной

словоформы (из-за морфологической омонимии их может быть несколько)

были собраны в одну запись. Кроме того, имеется возможность помимо

морфологической информации помещать любую другую, относящуюся к

данной словоформе лексическую или семантическую информацию. Это дает

возможность в перспективе сделать универсальный

    мор\-фо-се\-ман\-ти\-че\-ский словарь. Для объединения разнородной

информации в одной записи при формировании файла ММ осуществляется

специальный процесс сборки, когда из разных источников в~одну запись

собирается информация о~данной словоформе.

{\looseness=1



}



    Процесс сборки происходит в несколько стадий с помощью того же

блока загрузки, но работающего в режиме <<склеивания>> с последующей

выгрузкой промежуточного варианта МИ в текстовый вид.



%\vspace*{-6pt}



\section{ Некоторые параметры новой версии}



\vspace*{-2pt}



    Общий объем двух массивов (МА и МИ) для словаря в 90~тыс.\ основ

со\-став\-ля\-ет 120~МБ (в старой версии~--- менее 1~МБ).



    Время загрузки словаря~--- порядка 10~с.



    Скорость обработки текста русской деловой прозы увеличилась

приблизительно в 4~раза. При этом если учесть, что около 10\% слов

распознаются <<по аналогии>>, т.\,е.\ со старой скоростью, то скорость

обработки слов, имеющихся в словаре, увеличилась в~6~раз. Следует

отметить, что сам механизм доступа к~морфологической информации гораздо

более эффективен (по крайней мере в~десятки раз). Однако анализ показал,

что снижение эффективности происходит из-за того, что все процессы

<<постморфологии>>, занимающие достаточно много времени, остались от

старой версии.



\vspace*{-2pt}



\section{Возможности дальнейшего развития}



\vspace*{-2pt}



    В порядке развития данного проекта намечаются следующие работы:

    \begin{itemize}

\item перевод всей морфологической системы под ОС UNIX;

    \item создание новой версии для английской морфологии;

    \item перенесение некоторых вычислений из фазы <<постморфологии>>

на этап сборки информации при вводе~--- это позволит еще более повысить

эффективность анализа;

    \item поскольку информация в словаре имеет строго текстовый вид, то

появляется возможность реализации на этой основе других компонент

лингвистического процессора, например предметных терминологических

словарей.

    \end{itemize}



\vspace*{-2pt}



{\small\frenchspacing

{%\baselineskip=10.8pt



\begin{thebibliography}{99}



\vspace*{-2pt}



    \bibitem{1-som}

    \Au{Сегалович И., Маслов М.} Русский морфологический анализ и

синтез с генерацией моделей словоизменения для не описанных в словаре

слов~/\!\!/ Диалог'98: Тр. Междунар. семинара по компьютерной лингвистике

и ее приложениям.~--- Казань: Хэтер, 1998. Т.~2. С.~547--552.

\bibitem{2-som}

\Au{Коваленко А.} Вероятностный морфологический анализатор русского и

украинского языков~/\!\!/ Системный администратор, 2002. №\,1. {\sf

http://www.\linebreak keva.ru/stemka/stemka.html}.

    \bibitem{3-som}

    \Au{Сокирко А.\,В.} Морфологические модули на сайте {\sf

www.aot.ru}~/\!\!/  Компьютерная лингвистика и интеллектуальные

технологии: Тр. Междунар. конф. <<Диалог-2004>>.~--- М.: Наука, 2004.

C.~559--564.

    \bibitem{4-som}

    \Au{Сомин Н.\,В., Соловьева Н.\,С., Кузнецова Е.\,В., Шарнин~М.\,М.}

Система морфологического анализа: опыт эксплуатации и модификации~/\!\!/

Системы и средства информатики, 2005. Вып.~15. С.~20--30.

    \bibitem{5-som}

    \Au{Кузнецов И.\,П. Сомин Н.\,В.} Анг\-ло-рус\-ская сис\-те\-ма

извлечения знаний из потоков информации в среде Интернет~/\!\!/ Системы и

средства информатики, 2007. Вып.~17. С.~236--254.



    \bibitem{8-som} %6

    \Au{Кузнецов И.\,П., Мацкевич А.\,Г.}

    Се\-ман\-ти\-ко-ори\-ен\-ти\-ро\-ван\-ные сис\-те\-мы на основе баз

знаний.~--- М.: МТУСИ, 2007. 173~с.

    \bibitem{6-som} %7

    \Au{Сомин Н.\,В., Кузнецов И.\,П., Мацкевич А.\,Г., Николаев~В.\,Г.}

Методы и средства настройки мор\-фо-лек\-си\-че\-ско\-го анализатора на

предметную область~/\!\!/ Системы и~средства информатики,

2009. Вып.~19. С.~96--118.



    \bibitem{10-som} %8

    \Au{Кузнецов И.\,П., Сомин Н.\,В.} Особенности

    лек\-си\-ко-мор\-фо\-ло\-ги\-че\-ско\-го анализа при извлечении

информационных объектов и связей из текстов естественного языка~/\!\!/

Компьютерная лингвистика и интеллектуальные технологии: По мат-лам

ежегодной Междунар. конф. <<Диалог>>.~---  М.: РГГУ, 2010.

    Вып.~9(16). С.~254--264.

    \bibitem{7-som} %9

    \Au{Сомин Н.\,В., Кузнецов И.\,П., Николаев В.\,Г., Соловьева~Н.\,С.,

Мацкевич~А.\,Г.} Методы устранения неопределенностей блока

    лек\-си\-ко-мор\-фо\-ло\-ги\-че\-ско\-го анализа при извлечении знаний

из текстов естественного языка~/\!\!/ Системы и~средства информатики,

2011. Т.~21. №\,2. С.~96--114.



    \bibitem{9-som} %10

    \Au{Кузнецов И.\,П., Сомин Н.\,В., Козеренко~Е.\,Б.} Особенности

лек\-си\-ко-мор\-фо\-ло\-ги\-че\-ско\-го анализа в задачах извлечения

структур знаний из текстов естественного языка~/\!\!/ Искусственный

интеллект (журнал НАН Украины), 2011. Т.~3. С.~105--116.



    \bibitem{11-som}

    \Au{Кузнецов И.\,П., Сомин Н.\,В.} Выявление имплицитной

информации из текстов на естественном языке: проблемы и методы~/\!\!/

Информатика и её применения, 2012. Т.~6. Вып.~1. С.~49--58.

    \bibitem{12-som}

    \Au{Соловьева Н.\,С., Сомин Н.\,В.} ТЕРМИН-3~--- сис\-те\-ма

динамического гипертекста~/\!\!/ Сис\-те\-мы и средства информатики, 1995.

Вып.~7. С.~95--104.

    \bibitem{13-som}

    \Au{Дягилева А.\,В., Киселев С.\,Л., Сомин~Н.\,В.} Статистическая

модель рубрикации текстов на примере сообщений СМИ~/\!\!/

Дистанционное образование, 1998. №\,7. С.~16--21.

\bibitem{14-som}

\Au{Сомин Н.\,В., Соловьева Н.\,С.} Рубрицирование текстов как

информационная технология~/\!\!/ Системы и средства информатики, 2001. Вып.~11. С.~195--201.

    \bibitem{15-som}

    \Au{Сомин Н.\,В.} Каталоги специальных знаний в системе анализа

текстов на естественном языке~/\!\!/ Проблемы и методы информатики:

Тезисы докл. II~Научной сессии Института проблем информатики

Российской академии наук.~--- М.: ИПИ РАН, 2005. С.~128--131.

\end{thebibliography}

} }





\vspace*{-6pt}



\hfill{\small\textit{Поступила в редакцию 12.08.14}}



%\newpage





\vspace*{8pt}



\hrule



\vspace*{2pt}



\hrule



%\vspace*{-6pt}





\def\tit{USING HASH FUNCTION FOR~INCREASING SPEED OF~WORK OF~THE~SOFTWARE

FOR~MORPHOLOGICAL ANALYSIS OF~RUSSIAN TEXTS}



\def\titkol{Using hash function for~increasing speed of~work of~the~software

for~morphological analysis} % of~russian texts}



\def\aut{N.\,V.~Somin and M.\,M.~Sharnin}

\def\autkol{N.\,V.~Somin and M.\,M.~Sharnin}





\titel{\tit}{\aut}{\autkol}{\titkol}



\vspace*{-12pt}



\noindent

Institute of Informatics Problems, Russian Academy of Sciences,

44-2 Vavilov Str., Moscow 119333, Russian

Federation







\def\leftfootline{\small{\textbf{\thepage}

\hfill Sistemy i Sredstva Informatiki~---

Systems and Means of Informatics 2014 vol~24 no\,3}

}%

 \def\rightfootline{\small{Sistemy i Sredstva Informatiki~---

 Systems and Means of Informatics   2014   vol~24   no\,3

\hfill \textbf{\thepage}}}





\Abste{The paper considers the problem of increasing efficiency of morphological analysis of

Russian texts. The software system for morphological\linebreak\vspace*{-12pt}}



\Abstend{analysis is described, including the set of

morphological characteristics and the algorithms of work. The paper mentions the software

systems solving the problem of logiс-semantic analysis of natural language texts in which the

software system for morphological analysis found application. Features of the system are

discussed from the point of view of occupied memory and work speed. The way of storage of

morpholexical information using hash functions is suggested which provides high speed of access.

The difficulties arising during realization of such approach are discussed and possible solutions

are considered.  The paper describes the structure of information arrays of the new version and

the search algorithms realized in it. The paper also describes a subsystem for putting in and

updating morphological information.   Specific parameters of the new realization of the

software system for morphological analysis and information on speed of work acceleration in

comparison with the previous version are given.  The paper discusses opportunities of

development of the new version of the software system for morphological analysis and of

transferring the suggested approach to other components of the linguistic processor.}



    \KWE{morphological analysis; hash function; linguistic processor; logiс-semantic

analysis of natural language texts}





\DOI{10.14357\!/\!08696527140315}



%\vspace*{-24pt}



%\Ack

%\noindent





%\vspace*{-2pt}





\renewcommand{\bibname}{\protect\rmfamily References}

%\renewcommand{\bibname}{\large\protect\rm References}



{\small\frenchspacing

{%\baselineskip=10.8pt

\addcontentsline{toc}{section}{References}

\begin{thebibliography}{99}



\bibitem{1-som-1}

\Aue{Segalovich, I., and M.~Maslov}. 1998. Russkiy morfologicheskiy

analiz i sintez s~ge\-neratsiey modeley slovoizmeneniya dlya ne opisannykh

v slovare slov

[The Russian morphological analysis and synthesis with generation of models

of word change for the words not described in the dictionary].

\textit{Conference

Dialog'98 Proceedings}. Kazan': Hjeter. 2:547--552.



\bibitem{2-som-1}

\Aue{Kovalenko, A.} 2002. Veroyatnostnyy morfologicheskiy analizator russkogo i

ukrainskogo yazykov [Probabilistic morphological analyzer of the Russian and

Ukrainian languages]. \textit{Sistemnyy Administrator} [System Administrator] 1.

Available at:

{\sf http://www.keva.ru/stemka/stemka.html} (accessed August~12, 2014).



\bibitem{3-som-1}

\Aue{Sokirko, A.\,V.} 2004. Morfologicheskie moduli na sayte {\sf www.aot.ru}

[Morphological modules at the site {\sf www.aot.ru}].

\textit{Conference

Dialog'2004 Proceedings}.  Verkhnevolzhskiy. Available at:

{\sf http://www.aot.ru/docs/sokirko/Dialog2004.htm} (accessed August~12, 2014).



\bibitem{4-som-1}

\Aue{Somin, N.\,V., N.\,S. Solov'eva, E.\,V.~Kuznetsova, and M.\,M.~Sharnin}.

2005. Sistema

morfologicheskogo analiza: Opyt ekspluatatsii i modifikatsii [The System of

morphological analysis: Operating experience and modifications].

\textit{Sistemy i Sredstva Informatiki}~---

\textit{Systems and Means of Informatics} 15:20--30.



\bibitem{5-som-1}

\Aue{Kuznetsov, I.\,P., and N.\,V. Somin}. 2007. Anglo-russkaya sistema izvlecheniya

znaniy iz potokov informatsii v srede Internet [English-Russian languages system

of extraction of knowledge from Internet].

\textit{Sistemy i Sredstva Informatiki}~--- \textit{Systems

and Means of Informatics} 17:236--254.



\bibitem{8-som-1} %6

\Aue{Kuznetsov, I.\,P., and A.\,G.~Matskevich}. 2007.

\textit{Semantic-orientirovannye sistemy

na osnove baz znaniy} [Semantiko-focused systems on the basis of knowledge

bases]. Moscow: MTUCI. 173~p.



\bibitem{6-som-1} %7

\Aue{Somin, N.\,V., I.\,P. Kuznecov, A.\,G.~Matskevich, and V.\,G.~Nikolaev}.

2009. Metody

i sredstva nastroyki morfo-leksicheskogo analizatora na predmetnuyu oblast'

[Methods and means of setup of the morfolexical analyzer for subject domain].

\textit{Sistemy i Sredstva Informatiki}~---

\textit{Systems and Means of Informatics} 19:96--118.



\bibitem{10-som-1} %8

\Aue{Kuznetsov, I.\,P., and N.\,V.~Somin}. 2010.

Osobennosti leksiko-morfologicheskogo

analiza pri izvlechenii informatsionnykh ob"ektov i svyazey iz tekstov

estestvennogo yazyka [Features of the lexical-morphological analysis at extraction

of information objects and communications from natural language texts].

\textit{Komp'yuternaya Lingvistika i Intellektual'nye Tekhnologii. Po mat-lam

Mezhdunar. Konf. ``Dialog 2010''} [Computational Linguistics and

Intellectual Technologies. Conference (International) Dialog 2010 Proceedings].

Moscow: RGGU. 9(16):254--264.





\bibitem{7-som-1} %9

\Aue{Somin, N.\,V., I.\,P. Kuznetsov, V.\,G.~Nikolaev, N.\,S.~Solovyeva,

and A.\,G.~Matskevich.}

2011. Metody ustraneniya neopredelennostey bloka leksiko-morfologicheskogo

analiza pri izvlechenii znaniy iz tekstov estestvennogo yazyka [Methods to

eliminate uncertainty of lexical-morphological analysis

in process  of knowledge extraction from natural language texts]. \textit{Sistemy i Sredstva Informatiki}~---

\textit{Systems and Means of Informatics} 21(2):96--114.





\bibitem{9-som-1} %10

\Aue{Kuznetsov, I.\,P., N.\,V.~Somin, and E.\,B.~Kozerenko}. 2011.

Osobennosti leksiko-morfologicheskogo analiza v zadachakh izvlecheniya

struktur znaniy iz tekstov estestvennogo yazyka

[Features of the lexical-morphological analysis in problems

of extraction of structures of knowledge from natural language texts].

\textit{Iskusstvennyy Intellekt (NAN Ukrainy)}

[Artificial Intelligence (National

Academy of Sciences of Ukraine)] 3:105--116.





\bibitem{11-som-1}

\Aue{Kuznetsov, I.\,P., and N.\,V.~Somin}. 2012.

Vyyavlenie implitsitnoy informatsii iz

tekstov na estestvennom yazyke: Problemy i metody [Extraction of implicit

information from the natural language texts: Problems and methods].

\textit{Informatika i ee Primeneniya}~--- \textit{Inform. Appl}. 6(1):49--58.



\bibitem{12-som-1}

\Aue{Solovyeva, N.\,S., and N.\,V.~Somin}. 1995.

TERMIN-3~--- sistema dinamicheskogo

giperteksta [The system of  dynamic hypertext].

\textit{Sistemy i Sredstva Informatiki}~---

\textit{Systems and Means of Informatics}

 7:95--104.



\bibitem{13-som-1}

\Aue{Dyagileva, A.\,V., S.\,L.~Kiselev, and N.\,V.~Somin}. 1998.

Statisticheskaya model'

rubrikatsii tekstov na primere soobshcheniy SMI [The statistical model of a

rubrication of texts on mass media examples].

\textit{Distantsionnoe Obrazovanie} [Remote Education] 7:16--21.



\bibitem{14-som-1}

\Aue{Somin, N.\,V., and N.\,S.~Solovyeva}. 2001. Rubritsirovanie tekstov kak

informatsionnaya tekhnologiya [Rubrication of texts as information technology].

\textit{Sistemy i Sredstva Informatiki}~---

\textit{Systems and Means of Informatics} 11:195--201.



\bibitem{15-som-1}

\Aue{Somin, N.\,V.} 2005.

Katalogi spetsial'nykh znaniy v sisteme analiza tekstov na

estestvennom yazyke [Catalogs of special knowledge in system of the analysis of

texts in a~natural language].

\textit{Problemy i Metody Informatiki. II~Nauchnaya Sessiya

IPI RAN}.  Tezisy dokladov. [Problems and

Methods of Informatics. II~Scientific Session of IPI RAN. Abstracts]. Moscow: IPI RAN. 128--131.

\end{thebibliography}

} }





%\end{multicols}



%\vspace*{-12pt}



\hfill{\small\textit{Received August 12, 2014}}



%\vspace*{-24pt}



\Contr



\noindent

\textbf{Somin Nicolay V.} (b.\ 1947)~---

 Candidate of Science (PhD) in physics and mathematics, leading scientist, Institute of Informatics Problems, Russian Academy of Sciences,

44-2 Vavilov Str., Moscow 119333, Russian Federation;  somin@post.ru



\vspace*{3pt}



\noindent

\textbf{Sharnin Mihael M.} (b.\ 1959)~---

Candidate of Science (PhD) in technology, senior scientist, Institute of Informatics Problems, Russian Academy of Sciences,

44-2 Vavilov Str., Moscow 119333, Russian Federation;  mc@keywen.com





\label{end\stat}



\renewcommand{\bibname}{\protect\rm Литература}







